\documentclass[conference]{IEEEtran}

\usepackage{amsmath,amssymb}
\usepackage{booktabs}
\usepackage{graphicx}
\usepackage{multirow}
\usepackage{array}
\usepackage{xcolor}
\usepackage{url}

\newcommand{\todo}[1]{\textcolor{red}{[TODO: #1]}}
\newcommand{\placeholderfigure}[4]{%
  \begin{figure}[t]
    \centering
    \fbox{\begin{minipage}[c][#1][c]{0.94\columnwidth}
      \centering
      \textbf{Illustration placeholder}\\[3pt]
      \small #4
    \end{minipage}}
    \caption{#2}
    \label{#3}
  \end{figure}}

\title{A Bit-Exact Optimized Pipelined FP32 Approximate\\
Divider--Multiplier with Fixed Accuracy Levels}

\author{%
  \IEEEauthorblockN{\todo{Author names}}
  \IEEEauthorblockA{\todo{Affiliation, city, country, and email}}
}

\begin{document}
\maketitle

\begin{abstract}
Approximate floating-point arithmetic can trade bounded numerical error for
lower implementation cost in error-tolerant workloads. This paper presents a
pipelined FP32 approximate divider--multiplier based on a shared tangent-plane
datapath. Multiplication and division reuse the same multi-level correction
network, while division applies a small reciprocal-square coefficient selected
from a fixed-level lookup table. Five separately synthesizable approximation
levels are provided, and the datapath accepts one operation per cycle with a
seven-cycle latency. A bit-exact optimization fuses reciprocal lookup and
constant scaling, reducing synthesized cell area by 4.03--6.55\% and
vectorless post-synthesis power by 6.37--11.48\% without changing any tested
32-bit result. In a TSMC 65-nm flow, the shared unit reduces area by 3.79\%,
7.39\%, 9.10\%, and 10.43\% relative to separately synthesized divider and
multiplier blocks at levels L1--L4, respectively. All fixed-level designs meet
a 2-ns PrimeTime constraint with nonnegative setup and hold slack. A routed
runtime-configurable implementation further demonstrates physical-design
feasibility at 500 MHz.
\end{abstract}

\begin{IEEEkeywords}
approximate computing, floating-point arithmetic, divider, multiplier,
pipelining, ASIC, shift-add, lookup table
\end{IEEEkeywords}

\section{Introduction}
Floating-point (FP) division is substantially more expensive than simple
addition and commonly has higher latency and energy than multiplication.
Approximate computing addresses this cost by replacing exact arithmetic with
controlled approximations when an application can tolerate numerical error
\cite{han2013approx}. Existing approximate FP multipliers exploit
linearization, truncation, or logarithmic transformations
\cite{chen2020oam,imani2019approxl}. Approximate dividers additionally use
reciprocal approximation, iterative refinement, or piecewise fitting
\cite{saadat2019divider,wen2024pace}. However, instantiating an approximate
divider and multiplier independently duplicates sign, exponent, mantissa
partitioning, correction, and normalization hardware.

The OADM analytical formulation observes that the tangent-plane
approximations of multiplication and division have closely related correction
terms \cite{oadmtheory}. This work turns that formulation into a complete
pipelined FP32 RTL and ASIC evaluation. The implementation supports a shared
multiply/divide mode, elaboration-time fixed accuracy levels, and a fused
reciprocal scaling structure. The fixed-level variants avoid charging every
design for runtime level-selection hardware and enable a direct comparison
against separately synthesized arithmetic units.

The main contributions are:
\begin{itemize}
  \item a seven-cycle FP32 divider--multiplier that shares a multi-level
        tangent-plane correction datapath between both operations;
  \item five fixed-level RTL variants, L0--L4, whose unused correction and LUT
        branches are removed at elaboration and synthesis;
  \item a bit-exact fusion of the reciprocal-square LUT and constant scaler
        that reduces registers and arithmetic-selection logic without changing
        coefficients, precision, truncation, or latency;
  \item RTL accuracy characterization, TSMC 65-nm synthesis and PrimeTime PPA
        analysis, and a routed runtime-configurable implementation.
\end{itemize}

\placeholderfigure{1.52in}{Conceptual organization of the proposed FP32
divider--multiplier. Sign and exponent processing surround a shared mantissa
correction datapath and mode-selectable reciprocal scaling.}{fig:overview}{%
Draw FP32 unpacking on the left; shared correction plane in the center;
division-only reciprocal scale and mode mux near the output; normalization and
special-value handling on the right. Mark the seven pipeline stages.}

\section{Approximation Formulation}
\subsection{Floating-Point Decomposition}
For normalized single-precision operands $F_x$ and $F_y$, write
\begin{equation}
  F_x=(-1)^{s_x}x2^{E_x}, \qquad
  F_y=(-1)^{s_y}y2^{E_y}, \qquad x,y\in[1,2).
\end{equation}
The output sign is $s_x\oplus s_y$. Multiplication adds exponents and division
subtracts them; the expensive operation is therefore concentrated in the
mantissa result $xy$ or $x/y$.

\subsection{Tangent-Plane Approximations}
Consider a rectangular mantissa region
$[x_1,x_2]\times[y_1,y_2]$ with midpoint
\begin{equation}
  k_2=\frac{x_1+x_2}{2}, \qquad k_1=\frac{y_1+y_2}{2}.
\end{equation}
The tangent plane for multiplication is
\begin{equation}
  xy \approx z_M=-k_1k_2+k_1x+k_2y.
  \label{eq:mulplane}
\end{equation}
The corresponding divider approximation is
\begin{equation}
  \frac{x}{y}\approx z_D=
  \frac{k_1k_2+k_1x-k_2y}{k_1^2}.
  \label{eq:divplane}
\end{equation}
Thus, both modes share additions, signed corrections, midpoint constants, and
mantissa partition information. Division additionally scales the shared
numerator by $1/k_1^2$.

\subsection{Multi-Level Correction}
At level $n$, the mantissa domain is divided into $2^n$ intervals per axis.
Let $\hat{x}^{n}_{p(n)}$ and $\hat{y}^{n}_{q(n)}$ denote the midpoint of the
selected intervals. The unscaled divider numerator is
\begin{equation}
  w_n=\hat{x}^{n}_{p(n)}\hat{y}^{n}_{q(n)}
      +\hat{y}^{n}_{q(n)}x-\hat{x}^{n}_{p(n)}y,
\end{equation}
and it can be accumulated incrementally:
\begin{equation}
  w_n=w_0+\sum_{i=1}^{n}\Delta w_i,qquad
  w_0=2.25+1.5x-1.5y.
  \label{eq:divcorrection}
\end{equation}
The multiplier has an analogous expansion
\begin{equation}
  z_M^n=z_M^0+\sum_{i=1}^{n}\Delta z_i,qquad
  z_M^0=-2.25+1.5x+1.5y.
  \label{eq:mulcorrection}
\end{equation}
The OADM formulation relates the two correction sequences as
\begin{align}
  \Delta w_n &= -\Delta z_n
    + \left(2x\cdot\sigma_y(n)\right)\gg(n+1),
    \label{eq:shared_delta}\\
  w_0 &= -z_M^0+3x,
  \label{eq:shared_base}
\end{align}
where $\sigma_y(n)\in\{-1,+1\}$ is selected by the $n$th fractional bit of
$y$. Equations~\eqref{eq:shared_delta} and \eqref{eq:shared_base} motivate a
shared shift-add network instead of two complete arithmetic datapaths.

\placeholderfigure{1.32in}{Recursive partition of the normalized mantissa
domain for L0--L4. Each deeper level selects a finer midpoint and adds one
correction term.}{fig:partition}{%
Show the $[1,2)\times[1,2)$ square split into 1, 4, 16, 64, and 256 regions.
Highlight one nested region and label its midpoint $(\hat{x},\hat{y})$.}

\section{Hardware Architecture}
\subsection{Shared Datapath}
The top-level interface contains 32-bit operands $x$ and $y$, a clock, a
multiply/divide mode, and a 32-bit result. In multiply mode, the shared
correction sum directly enters normalization. In divide mode, the same sum is
multiplied by a reciprocal-square coefficient. Base terms and corrections are
implemented by shifts, conditional inversion, adders, and three 3:2 carry-save
compressors. No exact mantissa divider is instantiated.

The fixed-level wrappers make the selected approximation level an elaboration
parameter. L0 uses only the base plane. L1--L4 progressively retain one to four
correction terms. The reciprocal-square table has 1, 2, 4, 8, and 16 Q0.8
coefficients for L0--L4. Constant propagation removes unselected correction
levels and coefficient branches.

\subsection{Seven-Cycle Pipeline}
Table~\ref{tab:pipeline} summarizes the pipeline. The implementation accepts a
new operation every cycle; latency is seven cycles for both modes.

\begin{table}[t]
  \caption{Pipeline organization}
  \label{tab:pipeline}
  \centering
  \small
  \begin{tabular}{cl}
    \toprule
    Stage & Main operation \\
    \midrule
    S1 & FP classification, base plane, first correction \\
    S2 & Second correction and midpoint refinement \\
    S3 & Third/fourth corrections and level pruning \\
    S4 & Carry-save correction reduction and final addition \\
    S5 & Reciprocal constant scaling / multiply-mode bypass \\
    S6 & Mode alignment and mantissa selection \\
    S7 & Normalization, exponent adjustment, FP32 packing \\
    \bottomrule
  \end{tabular}
\end{table}

\placeholderfigure{1.48in}{Cycle-level pipeline and mode alignment. The
division scale and multiplication bypass converge before FP32 normalization.}
{fig:pipeline}{%
Draw S1--S7 boxes with registers between stages. Show exponent/sign/status
metadata flowing in parallel with the mantissa path.}

\subsection{Bit-Exact Fused Reciprocal Scaler}
The baseline divider stores an 8-bit coefficient selected by a separate LUT and
then applies a generic eight-term shift-add multiplier. This treats the
coefficient as an arbitrary runtime value even though each fixed-level design
can produce only a small finite set. The optimized structure combines the LUT
selection with full-width multiplication by the exact branch constant:
\begin{equation}
  r=\left(w_n C_{n,q}\right)\gg8,qquad
  C_{n,q}=\operatorname{Q0.8}\!\left(1/\hat{y}_{q(n)}^2\right).
\end{equation}
The constants and 38-bit two's-complement product semantics are unchanged.
Synthesis maps each branch to constant shift-add logic. Registering the complete
branch result replaces the wider partial-sum register array, while preserving
the original stage boundary and truncation point.

\subsection{FP32 Behavioral Scope}
The RTL preserves FP32 sign, exponent, and fraction fields and handles zero,
infinity, and NaN cases. Inputs with an all-zero exponent are treated as zero;
subnormal arithmetic is therefore flushed rather than gradually normalized.
The mantissa result is truncated during packing rather than rounded with full
IEEE round-to-nearest-even logic. These choices are held constant in all
baseline, fused, shared, and separate comparisons.

\section{Verification and Experimental Methodology}
\subsection{Functional and Numerical Verification}
RTL simulation was performed in ModelSim 2024.2. The numerical testbench
compared each approximate output against real-valued $xy$ or $x/y$ over 1007
finite normalized cases per operation and level. Metrics include mean relative
error distance (MRED), relative root-mean-square error (R-RMSE), and maximum
relative error.

To validate the fused scaler, the original and optimized scalers were compared
over 100,000 deterministic value/index vectors for all levels. A second
cycle-accurate regression recorded 20,000 cycles from all five complete FP32
fixed-level units, including both modes and zero, infinity, NaN, and exponent
boundary cases. Baseline and optimized traces were byte-identical, with the
same SHA-256 digest
\texttt{0b356af3b865e244...40379101cbe6403e1b}.

\subsection{ASIC Flow}
The fixed-level designs were synthesized with Synopsys Design Compiler
U-2022.12-SP7 using the TSMC 65-nm typical CCS library. DC used a 1.5-ns clock
to provide implementation guardband. Post-synthesis timing and vectorless power
were evaluated in PrimeTime U-2022.12-SP5 at a 2-ns published clock period.
Identical vectorless activity assumptions were used for every fixed-level
comparison: input static probability 0.5 and toggle rate 0.1 transitions/ns.
Consequently, these power values support relative comparisons but are not
workload-specific signoff power.

For physical feasibility, the runtime-configurable L0--L3 unit was placed,
clock-tree synthesized, and routed in Cadence Innovus 21.16. Extracted
post-route timing and time-based power were analyzed in PrimeTime. The
fixed-level variants have not yet been individually placed and routed; their
PPA tables therefore report cell area and post-synthesis power.

\section{Results}
\subsection{Approximation Accuracy}
Table~\ref{tab:accuracy} reports actual RTL simulation results. Each additional
level sharply reduces error in both modes. At L3, multiplication and division
R-RMSE are 0.0666\% and 0.1741\%, respectively.

\begin{table}[t]
  \caption{RTL numerical error on normalized finite inputs (\%)}
  \label{tab:accuracy}
  \centering
  \small
  \begin{tabular}{ccrrr}
    \toprule
    Mode & Level & MRED & R-RMSE & Max rel. \\
    \midrule
    \multirow{4}{*}{Mul}
      & L0 & 3.0629 & 4.5431 & 25.0000 \\
      & L1 & 0.7645 & 1.0762 & 6.2500 \\
      & L2 & 0.1971 & 0.2751 & 1.5625 \\
      & L3 & 0.0478 & 0.0666 & 0.3906 \\
    \midrule
    \multirow{4}{*}{Div}
      & L0 & 4.0541 & 6.3678 & 33.2030 \\
      & L1 & 1.0590 & 1.6228 & 8.3808 \\
      & L2 & 0.3985 & 0.5516 & 1.9979 \\
      & L3 & 0.1423 & 0.1741 & 0.5982 \\
    \bottomrule
  \end{tabular}
\end{table}

\placeholderfigure{1.42in}{Relative RMSE versus approximation level for
multiplication and division, preferably using a logarithmic vertical axis.}
{fig:accuracy}{%
Plot the R-RMSE values in Table~\ref{tab:accuracy}. Use separate lines for Mul
and Div, and label every data point.}

\subsection{Bit-Exact Optimization}
Table~\ref{tab:fused} compares the baseline fixed-level implementation against
the fused reciprocal scaler. Cell area falls by 4.03--6.55\%. The number of
sequential cells falls at every level because the generic partial-sum register
array is removed. PrimeTime vectorless power falls by up to 11.48\%. All levels
meet both setup and hold constraints.

\begin{table*}[t]
  \caption{Fixed-level bit-exact optimization in TSMC 65 nm}
  \label{tab:fused}
  \centering
  \small
  \begin{tabular}{crrrrrrrr}
    \toprule
    & \multicolumn{3}{c}{Cell area ($\mu\mathrm{m}^2$)}
    & \multicolumn{2}{c}{Sequential cells}
    & \multicolumn{2}{c}{PT power (mW)}
    & Setup/Hold (ns) \\
    \cmidrule(lr){2-4}\cmidrule(lr){5-6}\cmidrule(lr){7-8}
    Level & Baseline & Fused & Reduction & Baseline & Fused
          & Baseline & Fused & Fused \\
    \midrule
    L0 & 5602.32 & 5275.08 & 5.84\% & 325 & 304 & 2.7712 & 2.5623 & 0.51/0.04 \\
    L1 & 9025.20 & 8494.92 & 5.88\% & 470 & 406 & 4.1530 & 3.8883 & 0.51/0.04 \\
    L2 & 12053.52 & 11567.88 & 4.03\% & 617 & 509 & 5.4595 & 4.8426 & 0.51/0.04 \\
    L3 & 14380.56 & 13439.16 & 6.55\% & 699 & 583 & 6.2657 & 5.5509 & 0.50/0.01 \\
    L4 & 15804.36 & 14887.44 & 5.80\% & 725 & 611 & 6.6803 & 5.9134 & 0.50/0.01 \\
    \bottomrule
  \end{tabular}
\end{table*}

\placeholderfigure{1.38in}{Area and vectorless-power reduction from the
bit-exact fused reciprocal scaler.}{fig:fused}{%
Use grouped bars for baseline/fused area by level and a secondary panel for
baseline/fused power. Annotate percentage reductions.}

\subsection{Shared Versus Separate Hardware}
To isolate sharing benefit, each divider was resynthesized with divide mode tied
high. The multiplier baseline is the corresponding original fixed-level OAM
RTL from AM-Lib \cite{amlib}. All blocks use the same technology, timing budget,
drive, and output-load assumptions. The multiplier is combinational and has
less complete special-value behavior than the proposed pipelined unit; the
comparison is therefore conservative for the shared design but not an
equal-latency comparison.

As shown in Table~\ref{tab:sharing}, fixed control and pipeline overhead make
L0 1.09\% larger than the separate sum. Sharing becomes beneficial from L1 and
increases with correction depth, reaching 10.43\% at L4.

\begin{table*}[t]
  \caption{Shared divmul versus separate divider and OAM multiplier}
  \label{tab:sharing}
  \centering
  \small
  \begin{tabular}{crrrrrr}
    \toprule
    Level & Div only & Mul only & Separate sum & Shared divmul
          & Area saved ($\mu\mathrm{m}^2$) & Saving \\
    \midrule
    L0 & 3897.36 & 1320.84 & 5218.20 & 5275.08 & $-56.88$ & $-1.09\%$ \\
    L1 & 6667.20 & 2162.52 & 8829.72 & 8494.92 & 334.80 & 3.79\% \\
    L2 & 9477.00 & 3014.28 & 12491.28 & 11567.88 & 923.40 & 7.39\% \\
    L3 & 11056.32 & 3727.44 & 14783.76 & 13439.16 & 1344.60 & 9.10\% \\
    L4 & 12256.56 & 4364.64 & 16621.20 & 14887.44 & 1733.76 & 10.43\% \\
    \bottomrule
  \end{tabular}
\end{table*}

\placeholderfigure{1.40in}{Area of separate divider plus multiplier versus the
shared divider--multiplier at each fixed approximation level.}{fig:sharing}{%
Use paired bars for Separate Sum and Shared DivMul. Add a line or labels for
the shared saving percentage.}

\subsection{Physical-Design Demonstration}
The runtime-configurable unit was routed with 12,230 placed instances including
physical-only cells. Its core area is $35{,}640\,\mu\mathrm{m}^2$, with
$20{,}356.26\,\mu\mathrm{m}^2$ of functional standard-cell area after
excluding physical-only cells. Innovus reports worst setup and hold slack of
+0.072 ns and +0.029 ns across the reported post-route views. In-tool geometry
and antenna checks report no violations. PrimeTime reports 5.3013 mW
time-based power at a 2-ns clock, corresponding to 10.6026 pJ per cycle. This
layout predates the fixed-level fused scaler and is presented only as evidence
that the shared pipelined architecture is routable and timing-clean; a new
fixed-level APR comparison remains future work.

\placeholderfigure{1.65in}{Routed layout of the runtime-configurable OADM
divider--multiplier in TSMC 65 nm.}{fig:layout}{%
Insert a clean Innovus/Virtuoso layout screenshot with the core boundary,
power grid, standard-cell region, and dimensions annotated.}

\section{Discussion}
The results expose two forms of efficiency. First, specializing reciprocal
scaling removes generic hardware while preserving every tested output bit.
Second, sharing becomes increasingly valuable as approximation depth grows,
because the correction network accounts for a larger fraction of total area.
The L0 result also shows that sharing is not free: mode alignment, pipeline
registers, and complete FP handling can dominate a minimal base approximation.

Several limitations define the scope of the current claims. Power for the
fixed-level designs is vectorless and should be replaced by workload-derived
SAIF activity for final energy claims. The OAM multiplier comparison is not
latency- or special-case-equivalent. L4 has been verified against the baseline
L4 RTL but has not yet been included in the real-valued accuracy campaign.
Finally, fixed-level APR, extracted multi-corner signoff, and Calibre LVS/DRC
results are not yet available for the fused variants.

\section{Conclusion}
This paper presented a pipelined FP32 approximate divider--multiplier that
shares tangent-plane base and correction hardware across both operations. A
fixed-level implementation and bit-exact fused reciprocal scaler reduce area
and vectorless power without changing tested arithmetic outputs or latency.
At L4, the fused unit occupies $14{,}887.44\,\mu\mathrm{m}^2$, consumes an
estimated 5.9134 mW post-synthesis, and is 10.43\% smaller than the sum of a
separate divider and OAM multiplier. The measured accuracy progression and
positive timing slack demonstrate a practical quality--cost tradeoff suitable
for further workload-driven and post-layout evaluation.

\begin{thebibliography}{99}

\bibitem{han2013approx}
J. Han and M. Orshansky, ``Approximate computing: An emerging paradigm for
energy-efficient design,'' in \emph{Proc. IEEE European Test Symposium}, 2013,
pp. 1--6.

\bibitem{chen2020oam}
C. Chen, S. Yang, W. Qian, M. Imani, X. Yin, and C. Zhuo, ``Optimally
approximated and unbiased floating-point multiplier with runtime
configurability,'' in \emph{Proc. IEEE/ACM ICCAD}, 2020, pp. 1--9.

\bibitem{imani2019approxl}
M. Imani \emph{et al.}, ``ApproxLP: Approximate multiplication with
linearization and iterative error control,'' in \emph{Proc. ACM/IEEE DAC},
2019, pp. 1--6.

\bibitem{saadat2019divider}
H. Saadat \emph{et al.}, ``Approximate integer and floating-point dividers with
near-zero error bias,'' in \emph{Proc. ACM/IEEE DAC}, 2019, pp. 1--6.

\bibitem{wen2024pace}
C. Wen \emph{et al.}, ``PACE: A piece-wise approximate and configurable
floating-point divider for energy-efficient computing,'' in
\emph{Proc. DATE}, 2024.

\bibitem{oadmtheory}
\todo{Authors}, ``OADM: Optimally approximate floating-point
divider--multiplier with runtime configurability,'' working manuscript,
\todo{year}.

\bibitem{ieee754}
\emph{IEEE Standard for Floating-Point Arithmetic}, IEEE Std. 754-2019, 2019.

\bibitem{amlib}
Y. Wu \emph{et al.}, ``AM-Lib: A library of approximate multipliers,'' OAM RTL,
commit \texttt{cc3864b}. [Online]. Available:
\url{https://github.com/skycrapers/AM-Lib/tree/main/OAM}

\end{thebibliography}

\end{document}
